\capitulo{7}{Conclusiones}
Tras completar el desarrollo, despliegue y validación de la herramienta web de recomendación de SAAC para pacientes con ELA, se presentan las principales conclusiones del proyecto divididas en los hitos alcanzados.

\begin{itemize}
    	\item \textbf{Cumplimiento de objetivos:} Se ha logrado diseñar e implementar un prototipo web funcional y accesible que centraliza las soluciones SAAC del mercado y automatiza su recomendación. La lógica del cuestionario traduce con éxito variables biomédicas y de 				accesibilidad en propuestas tecnológicas concretas, lo que optimiza un proceso que hasta ahora los profesionales realizaban de forma manual y dispersa.
    	\item \textbf{Validación en entornos reales:} Las pruebas con pacientes reales de ELACyL demostraron la utilidad práctica del sistema. Obtener una nota de adecuación de 6/10 en las primeras pruebas reales valida que el recomendador ofrece una guía realista, sirviendo en dos casos 			muy diferentes, tanto de respuesta inmediata para fases avanzadas como de herramienta de planificación preventiva en las fases iniciales.
    	\item \textbf{Impacto formativo e interdisciplinar:} Desde el punto de vista del aprendizaje, el proyecto ha permitido aplicar principios de Ingeniería de la Salud a un problema sociosanitario real. Se ha logrado unificar el lenguaje clínico (escalas de autonomía, fatiga o afectación del 		habla) con la lógica computacional (estructuras de bases de datos relacionales y condiciones lógicas de filtrado).
\end{itemize}

\section{Aspectos relevantes}
Esta sección resume la experiencia práctica acumulada durante las fases del proyecto, detallando las decisiones de diseño no triviales y los caminos descartados tras obtener resultados negativos.

\subsection{El diseño del modelo de datos y el efecto embudo}
Uno de los mayores retos del proyecto fue estructurar la base de datos de manera que reflejara fielmente el catálogo SAAC. Inicialmente, se planteó un modelo de emparejamiento demasiado rígido, donde cada perfil de usuario se conectaba directamente con un único kit cerrado de software y hardware. Sin embargo, las primeras simulaciones técnicas dieron resultados negativos: la ELA es una enfermedad muy heterogénea y este enfoque provocaba que muchos perfiles se quedaran sin ninguna recomendación válida al no cumplir con el 100\% de los requisitos estrictos de un kit.

Para solucionar esto, se optó por cambiar la estrategia hacia un modelo de filtrado por exclusión (efecto embudo). La base de datos se dividió en entidades independientes (software, hardware, periféricos, bancos de voz). El cuestionario web no busca el sistema ideal sino que va descartando lo que el paciente no puede usar. 

Por ejemplo, si el usuario marca Interacción: ojos, el sistema elimina directamente los ratones y pantallas táctiles, pero mantiene abiertos todos los \textit{Eye trackers} y softwares compatibles de la base de datos. Esto permite que el sistema trate de ofrecer alternativas viables y escalables.

\subsection{Justificación de la plataforma}
La elección de desarrollar una aplicación web (en lugar de una aplicación de escritorio o móvil) se justificó por criterios de accesibilidad y despliegue. Un entorno web garantiza que los profesionales de las asociaciones puedan acceder al recomendador desde cualquier dispositivo (Windows, Mac, tablet) sin necesidad de instalaciones complejas, únicamente entrando en un enlace. Además, facilita la actualización en tiempo real de la base de datos a medida que surjan nuevos dispositivos o aplicaciones de comunicación en el mercado, asegurando la sostenibilidad del proyecto a largo plazo.

